% Tapis de Quillen, with Quillen's offprint and preprint (cote 162-5, typed part): archive cote 162-5, site page 48 (PDF page 49).
% typescript. Draft compared: none.
% AI-transcribed by Claude Opus 5.5 (claude-opus-5-5, adaptive thinking, effort high), Message Batches API, batch pilotnext (msgbatch_01A9sdu8SYq23bEtqX9tUncQ), 2026-11-13; not yet checked by eye.
\scanpage{162-5/48}
Le théorème de Thom est alors qu'un $x \hand{$\in$} H^{\cdot}(X)$ est dans S(X) sss il est
\xout{xxxxxx}\inserted[typed]{fixe.} par G, i.e. satisfait $A^{+}. x = 0$. (Pourquoi cette condition est-el% NOTE: right edge of the page cut, end of line not visible
le déjà nécessaire ?) Une autre façon de le voir est de dire que le
foncteur $H^{\cdot}$ considéré comme allant de la catégorie $\underline{S}$ précédente dans
la catégorie des \xout{\unsure{vectoriels}} représentations de G dans des vectoriels
sur $\underline{F}_2$ , est pleinement fidèle. \xout{\unsure{Comme G est}} Quand on prend l'enveloppe
pseudo-abélienne, on trouve une sous-catégorie pleine de la catégorie des
représentations de G, mais il est douteux qu'on trouve une catégorie abé-%
lienne, voire la catégorie de toutes les représentations de G de dimensio% NOTE: right edge cut
finie sur $\underline{F}_2$ . (Quelle est la structure de G, qui est prolisse; d'après
\margin{oui, car \unsure{ga-} \illegible{3} \illegible{3} \illegible{2} \deleted{\illegible{1}} \illegible{4} \illegible{3} \unsure{fini stable par} $\triangleright$}Quillen, il serait pro-unipotent ....). Ici G bien sûr joue le rôle d'un
groupe de Galois motivique, mais contrairement à la situation familière,
il semble qu'il soit fortement non réductif (et d'ailleurs le corps de
définition dans tout ceci est $\underline{F}_2$ , et non $\underline{Q}$ ! ).

Quillent donne une généralisation de la construction d'une caté-%
gorie universelle $\underline{B}$ au cas où on travaille avec des variétés pas nécessai% NOTE: right edge cut
rement propres (mais f. n'étant supposé défini que pour f propre), et
quand on introduit une condition d'orientation \inserted[typed]{stable} sur les \uline{morphismes}
envisagés , via un morphisme $H \hand{$\to$} BO(\infty )$, où H est un H-espace (commuta-%
tif, associatif et unitaire ?): une orientation de $f:X \hand{$\to$} Y$ sera un
relèvement homotopique de $X \hand{$\to$} BO$ défini par \xout{fxx(Tx} $T_X - f^{\cdot}(T_Y)$ en $X \hand{$\to$} H$.
les éléments de H(X,Y) seront définis par des diagrammes de morphismes
orientés $Z \hand{$\to$} X$, $Z \hand{$\leftarrow$} Y$, le deuxième étant propre, modulo bords de tels.
En tant qu'ensemble, il se définit aussi en termes de la vari té XxY
munie simplement de la famille des supports formée des ensembles fermés
propres sur Y ..... \hand{Les $B^{H}_{m}(X)$ (X compact disons) peuvent}
\hand{se décrire homotopiquement à la Thom, comme les groupes d'homotopie}
\hand{stables d'espaces \uline{Hom}(X, \unsure{$MH(N-m)$})} \hand{-----.}

\mspage{\hand{48}}
